% from main.tex
%%%%%%%%%%  scan idx 626  |  folio 613  %%%%%%%%%%
En
effet ces conditions sont suffisantes car déjà 4) et 5) entraînent l'effectivité. Réciproquement soit \(M\) inversible sur \(Y\) tel qu'il existe un isomorphisme \(f^*M\simeq L\) compatible avec la donnée de recollement. On peut recouvrir \(Y\) par des ouverts \(V\), au-dessus desquels \(M\) est isomorphe à \(\mathcal O_Y\)~; donc la donnée de descente sur \(L\) restreindre au-dessus de \(V=f^{-1}(V)\) est isomorphe à la donnée de descente triviale sur \(\mathcal O_X|V\). Enfin les conditions 1)-5) sont de nature locale sur \(Y\) et elles sont vérifiées d'après les hypothèses de 4.1 pour \(L=\mathcal O_X\)~; elles sont donc vérifiées pour \(L\).

Il reste donc à représenter successivement chacune des conditions 1) à 5) par des monomorphismes de présentation finie. Pour 1) et 2) on applique le

\textbf{Lemme 4.4.} Soient \(S\) un schéma, \(\lambda:X\to S\) et \(\mu:Y\to S\) des morphismes propres de présentation finie, \(f:X\to Y\) un \(S\)-morphisme. Soit \(F\) un faisceau sur \(X\) de présentation finie et plat sur \(S\) et, pour un entier \(n\), soit \(S_n\) le sous-foncteur de \(S\) qui rend \(R^if_*(F)\) plat sur \(S\) pour tout \(i\ge n\) (précisément, \(S_n(T)=\{\varnothing\}\) si \(R^1(f_T)_*(F_T)\) est plat pour tout \(i\ge n\), \(S_n(T)=\varnothing\) sinon). Alors \(S_n\hookrightarrow S\) est représentable par un monomorphisme surjectif de présentation finie, et alors pour \(i\ge n\) les \(R^i(f_{S_n})_*(F_{S_n})\) commutent aux changements de base \(T\to S_n\). De plus \(S_n\) coïncide avec le foncteur \(S'_n\) qui rend exact les foncteurs
\[
  M\longmapsto R^i(f_V)_*(F_V\otimes\lambda_V^*M)\qquad \text{pour } i\ge n
\]
où \(M\) parcourt les \(\mathcal O_V\)-modules quasi-cohérents pour tout ouvert affine \(V\) de \(S\).

En effet l'assertion est locale sur \(S\)~; on peut donc supposer \(S\) quasi-compact. Alors la dimension relative de \(X/Y\) est bornée, soit par \(r\), et pour \(n>r\), \(R^nf_*(F_V\otimes\lambda_V^*M)=0\). L'assertion est donc vraie
